\chapter{Literature Review}
\label{chap:literature_review}

This chapter reviews the literature through the logic of the dissertation rather than as a loose catalogue of related topics. The thesis contains two connected study focuses and four model blocks. Study~1 examines CSR-driven manufacturer--retailer coordination in an EV battery recycling channel. Its first model block compares operational channel regimes: Vertical Nash, Manufacturer Stackelberg, Retailer Stackelberg, and an integrated benchmark. Its second model block uses those operating-regime payoffs in a stage-0 meta-game that endogenizes whether firms coordinate or compete for leadership. Study~2 then extends the problem to a multi-retailer reverse-logistics setting. Its third model block formulates the EV battery recycling network as a MILP facility-location and flow problem, while its fourth model block allocates the resulting network cost using cooperative game theory and the Shapley value.

The review therefore asks a targeted question: what parts of the existing literature support these four blocks, and where do gaps remain? Table~\ref{tab:lit_synthesis} summarizes the organizing logic. The expanded paper-level coverage matrix is maintained separately in `literatures_all/literature_list.csv`.

\begin{table}[htbp]
\centering
\caption{Literature Streams, Model Blocks, and Thesis Positioning}
\label{tab:lit_synthesis}
\resizebox{\textwidth}{!}{%
\begin{tabular}{p{0.20\textwidth}p{0.22\textwidth}p{0.23\textwidth}p{0.27\textwidth}}
\toprule
\textbf{Dissertation block} & \textbf{Representative literature} & \textbf{Coverage in prior work} & \textbf{Remaining gap addressed here} \\
\midrule
Study~1 Block~1: operational bilateral CSR channel game & \cite{ZhangLiang2023}, \cite{Khosroshahi2019}, \cite{Ferrara2017}, \cite{Narang2024}, \cite{Zeng2024} & Two-echelon pricing, CSR effort, recycling investment, Stackelberg timing, and government incentives & Prior models often impose the operating regime and do not connect CSR-driven recycling maturity to a later endogenous regime-choice layer. \\
Study~1 Block~2: stage-0 meta-game equilibrium analysis & \cite{Buccella2024}, \cite{FantiBuccella2017}, \cite{ZhengB2022} & CSR commitment, network externalities, cooperation-mode choice, and strategic pre-commitment & Existing work rarely treats coordination versus leadership competition as a separate pre-operational equilibrium problem in EV battery recycling. \\
Study~2 Block~3: EV battery reverse-logistics network design & \cite{bafadal2024proposed}, \cite{tadaros2022location}, \cite{wang2025planning}, \cite{ma2024unlocking}, \cite{HoarauLorang2022} & Facility location, EPR obligations, battery return flows, echelon utilization, and MILP-based network design & Network models usually optimize physical flows centrally and leave stakeholder cost responsibility outside the optimization. \\
Study~2 Block~4: cooperative cost sharing & \cite{lozano2013cooperative}, \cite{meng2023profit}, \cite{luo2022core}, \cite{Xiao2024}, \cite{wang2023collaborative} & Shapley allocation, horizontal cooperation, constrained coalitions, and fairness/stability concepts & Cooperative-game models often lack an EV-battery-specific characteristic cost function generated by a reverse-logistics MILP. \\
\bottomrule
\end{tabular}}
\end{table}

\section{Study Focus 1: CSR-Driven Bilateral Coordination}
\label{sec:lit_study1}

Study~1 begins from a bilateral manufacturer--retailer channel. This structure is intentionally stylized: it isolates vertical power, price setting, and CSR-driven recycling maturity before the dissertation moves to the multi-retailer network in Chapter~\ref{chap:model2}. The relevant literature is not only EV-battery-specific; it also includes CSR game-theoretic models that clarify how social responsibility changes pricing, effort, and commitment.

\subsection{Model Block 1: Operational Channel Games}
\label{sec:lit_block1}

Operational channel models typically assume that a manufacturer and one or more downstream firms choose prices, effort, recycling investment, or channel roles under an imposed timing structure. In EV battery recycling, this logic appears in studies of closed-loop supply chains under government subsidies, reward--penalty policies, deposit-return systems, and echelon utilization. \cite{Gu2021} model secondary use under subsidy and show that the value of returned batteries depends on whether second-life channels are included. \cite{ZhangW2024} extend this logic by linking government subsidies and echelon utilization to optimal EV battery closed-loop supply chain decisions. \cite{Narang2024} compare recycling model selection under Stackelberg interaction, carbon cap-and-trade, and reward--penalty policy. \cite{Zeng2024} further show that different power structures under government subsidy generate different recycling and profit outcomes.

The CSR literature adds the behavioral and demand-side motivation for treating recycling maturity as more than a technical investment. \cite{ZhangLiang2023} introduce CSR and fairness concerns in a NEV battery closed-loop supply chain with third-party recycling, showing that social responsibility and fairness can alter coordination outcomes. More general CSR supply-chain game models reach a similar conclusion. \cite{Ferrara2017} study dynamic Stackelberg CSR allocation, while \cite{Khosroshahi2019} connect CSR, transparency-dependent demand, and pricing decisions. \cite{Mahdiraji2023} frame social responsibility as a factor that changes supply-chain profitability and contract performance.

These models motivate the operational part of Chapter~\ref{chap:model1}. The chapter keeps the bilateral structure but makes the recycling maturity variable explicitly CSR-driven: the retailer's investment in take-back, sorting, traceability, and recycling credibility raises demand because consumers and regulators value responsible battery end-of-life management. The VN, MS, RS, and integrated regimes therefore represent alternative operating games over price and CSR-driven recycling maturity.

\subsection{Model Block 2: Stage-0 Meta-Game and Endogenous Regime Choice}
\label{sec:lit_block2}

The second model block of Study~1 is not another operating regime. It is the stage-0 meta-game that precedes the VN/MS/RS/IN operating games. This distinction matters because many supply-chain models assume a regime first and then solve decisions inside that regime. In practice, however, the regime itself is often a strategic choice: the manufacturer and retailer may coordinate, attempt to lead, or avoid coordination because each fears being exploited by the other.

CSR game-theoretic work helps motivate this pre-operational layer. \cite{Buccella2024} show that CSR commitment can interact with network externalities and firm strategy. \cite{FantiBuccella2017} similarly treat CSR adoption in a game-theoretic context, making social responsibility a strategic object rather than an exogenous moral preference. In closed-loop supply chains, \cite{ZhengB2022} compare alliance and cost-sharing modes, showing that cooperation design itself can determine whether recycling coordination is attractive.

The gap is that these studies usually do not embed the coordination choice in an EV battery recycling model where operational payoffs have already been derived under VN, MS, RS, and integration. Chapter~\ref{chap:model1} addresses this by feeding the operating-regime payoffs into a stage-0 normal-form game. This allows the chapter to ask why coordination may fail even when integration is jointly efficient, and how fixed coordination costs or payoff-sharing rules affect equilibrium selection.

\section{Study Focus 2: Multi-Retailer Reverse Logistics and Cost Sharing}
\label{sec:lit_study2}

Study~2 extends the bilateral channel logic into a more operationally realistic reverse channel. Once multiple retailers collect batteries across a region, recycling is no longer only a pricing and effort problem. It becomes a physical network-design problem and a cost-allocation problem. The manufacturer may bear EPR responsibility, but retailers provide the geographically dispersed return points and battery volumes.

\subsection{Model Block 3: Reverse-Logistics Network Design}
\label{sec:lit_block3}

The EV battery recycling literature emphasizes that end-of-life batteries require collection, sorting, consolidation, transportation, and processing. Regulatory work explains why these flows matter. \cite{HoarauLorang2022} assess European battery recycling regulation, while \cite{yang2025bridging} review EPR implementation for power battery recycling in China. Technical reviews such as \cite{Harper2019}, \cite{tembo2024lithium}, and \cite{rezaei2025review} explain why recycling is environmentally and economically important, but they do not resolve how firms should share the reverse-channel burden.

Network models address the physical design problem. \cite{tadaros2022location} formulate facility-location and reverse-logistics design for lithium-ion batteries in Sweden. \cite{bafadal2024proposed} propose a MILP reverse-logistics network for end-of-life EV battery management in the Jakarta Greater Area. \cite{wang2025planning} add robust echelon utilization planning, and \cite{ma2024unlocking} develop an urban EV battery recycling optimization model. These studies establish the need for multi-echelon network design, but they generally treat the optimized network as if a single planner can implement and finance it.

This is the starting point for Chapter~\ref{chap:model2}. The chapter keeps the network-design strength of this literature, but it makes coalition-dependent costs explicit. The MILP is not only a planning model; it becomes the engine that generates the characteristic cost function used by the cooperative game.

\subsection{Model Block 4: Cooperative Game Theory and Cost Allocation}
\label{sec:lit_block4}

Cooperative game theory provides the allocation logic needed once a shared network is created. In logistics, \cite{lozano2013cooperative} use cooperative game theory to allocate the benefits of horizontal cooperation. \cite{wang2023collaborative} connect operational collaboration and Shapley allocation in carrier distribution. In broader supply-chain settings, \cite{meng2023profit} study profit allocation in a four-echelon supply chain, and \cite{luo2022core} review the core, Shapley value, nucleolus, and Nash bargaining solution in operations management.

EV battery recycling introduces a specific version of this allocation problem. \cite{Xiao2024} study cooperative EV battery recycling with blockchain technology and Shapley-based allocation, while \cite{ZhengB2022} compare alliance and cost-sharing modes in a closed-loop supply chain. These papers show that cooperation can improve system outcomes, but they do not fully integrate a facility-location MILP with a coalition-specific cost game for a multi-retailer EV battery network.

Chapter~\ref{chap:model2} fills that gap by linking the physical network and cooperative game directly. For each coalition, the reverse-logistics MILP determines the minimized network cost. The Shapley value then allocates the grand-coalition cost by each player's marginal contribution. This creates a bridge between network efficiency, budget balance, allocation fairness, and coalition stability.

\section{Synthesis and Research Gaps}
\label{sec:lit_gap}

Across these streams, the literature provides strong components but leaves important integration gaps. First, CSR and EV battery recycling models often study operational decisions under an assumed timing structure. They rarely make the coordination regime itself an endogenous stage-0 equilibrium object. Chapter~\ref{chap:model1} responds by treating VN, MS, RS, and integration as operating-regime payoffs that feed a pre-operational meta-game.

Second, EV battery reverse-logistics models are strong on facility location and material flows, but they often assume centralized implementation. Chapter~\ref{chap:model2} keeps the MILP network-design logic while asking how the resulting cost should be shared among independent retailers and the manufacturer.

Third, cooperative cost-sharing models provide fairness concepts, but many are not grounded in EV-battery-specific cost structures. Chapter~\ref{chap:model2} uses the MILP-generated characteristic cost function to make Shapley allocation operationally meaningful in a multi-retailer EV battery recycling network.

Together, the two studies therefore address a linked research gap: EV battery recycling supply chains need both CSR-driven strategic coordination and fair reverse-logistics cost sharing. Study~1 explains how vertical power and stage-0 coordination shape CSR-driven recycling maturity in a bilateral channel. Study~2 explains how that coordination problem scales when recycling becomes a shared network whose costs must be allocated across multiple stakeholders.
